\chapter{Methodology}
\label{ch:methodology}

\section{Research Questions and Hypotheses}
\label{sec:methodology-rqs}

\begin{itemize}
  \item \textbf{RQ1.} Does sentence-level factuality feedback aggregated via GCA
    produce a more learnable pairwise preference signal than holistic
    summary-level feedback, as measured by reward-model pairwise validation
    accuracy?

  \item \textbf{RQ2.} How do evaluator configuration and sentence-score
    aggregation strategy affect the relative learnability of holistic and
    GCA-derived preferences?

  \item \textbf{RQ3.} How stable are the observed differences across repeated
    reward-model training runs and dataset sizes?

  \item \textbf{RQ4.} Does the learnability advantage measured in RQ1 also
    hold when the trained reward models are scored against an independent,
    held-out ground truth constructed from the same factuality judge, rather
    than against the models' own training-style labels, and does that hold as
    the training-set size grows?
\end{itemize}

RQ1--RQ3 are addressed by the experiments in Chapter~\ref{ch:results} and
answered in Chapter~\ref{ch:conclusion}. RQ2 is addressed by an exploratory
sweep rather than a controlled ablation, so it supports statements about
association between configuration and outcome, not about causal mechanism
(Section~\ref{sec:methodology-exploratory}). RQ4 was added after the
$n=1{,}000$/$5{,}000$/$10{,}000$ campaign, in direct response to a limitation
that campaign left open: pairwise validation accuracy is measured against
labels supplied by the same evaluator that constructed the preferences, so a
reward model could in principle learn to recover those labels' idiosyncrasies
rather than genuine factuality signal (Section~\ref{sec:methodology-scope}).
Section~\ref{sec:methodology-ground-truth} describes the independent
evaluation constructed to address this, and
Section~\ref{sec:results-ground-truth} reports its results.

Three working hypotheses guided the design. H1: sentence-level aggregated
preferences are more learnable than holistic preferences, because they localize
the factuality signal within a long output. H2: the evaluator's configuration
affects whether any such difference is observable. H3: effects observed at one
dataset scale need not persist at another. Their status is assessed in
Chapter~\ref{ch:conclusion}.

\section{Study Design}
\label{sec:methodology-design}

The comparison is between two procedures for converting factuality scores into
pairwise preference labels. For every source article, two candidate summaries are
generated once and reused by both conditions:

\begin{description}
  \item[Holistic.] The complete candidate summary is scored against the source
    article, producing one score per candidate. The two scores are compared to
    yield a pairwise preference.
  \item[GCA.] The candidate summary is segmented into sentences; each sentence is
    scored separately against the source article; the sentence scores are
    aggregated into one summary-level score. The two aggregated scores are
    compared to yield a pairwise preference.
\end{description}

The following are held identical across conditions: the source articles; the two
candidate summaries and the decoding settings that produced them; the evaluator
family and checkpoint; the reward-model architecture; and the reward-model
training and evaluation procedure. The manipulated variable is the granularity at
which the evaluator is queried, together with the aggregation step that GCA
requires and holistic scoring does not.
Table~\ref{tab:methodology-summary} lists the final configuration.

\begin{table}[htbp]
\centering
\caption[Final experimental configuration]{Final experimental configuration.}
\label{tab:methodology-summary}
\begin{tabularx}{\textwidth}{lX}
\toprule
Factor & Value \\
\midrule
Task & Single-document abstractive summarization (news) \\
Dataset & CNN/DailyMail; runs at $n=1{,}000$, $5{,}000$, and $10{,}000$
  (subset seed 200) \\
Candidate generation & Mistral-7B-Instruct-v0.3, two-temperature sampling
  ($T=0.7$ / $T=1.0$) \\
Evaluator & \texttt{yzha/AlignScore}, mode \texttt{nli} \\
Manipulated variable & Feedback granularity: holistic vs.\ sentence-level $+$ GCA \\
GCA aggregation & Simple mean ($\alpha = 0.0$) \\
Margin & 0 (all pairs used) \\
Reward model & Bradley--Terry, \url{FacebookAI/roberta-base} backbone,
  mean-pool $+$ linear scalar head \\
RM training & epochs = 5, lr = $2\times10^{-5}$, batch size = 8,
  max length = 512, 5-fold cross-validation \\
Dependent variable & Reward-model pairwise validation accuracy \\
\bottomrule
\end{tabularx}
\end{table}

\section{Scope of the Evaluation}
\label{sec:methodology-scope}

The dependent variable throughout this thesis is the pairwise validation accuracy
of a Bradley--Terry reward model trained on a given preference set. Within this
thesis, \emph{learnability} denotes exactly this quantity: the extent to which the
preference relation encoded in a dataset can be recovered by the chosen
reward-model architecture under a fixed training procedure.

Because this term is easily over-read, it is worth stating precisely what a higher
accuracy does and does not establish. A higher value indicates that the generated
preference relation is more consistently recoverable by the selected reward-model
architecture. It does not by itself establish any of the following:

\begin{itemize}
  \item that the preference labels agree better with human factuality judgments;
  \item that the labels are objectively more accurate as factuality annotations;
  \item that a summarization model trained on those preferences produces more
    factually consistent summaries;
  \item that GCA is generally preferable to holistic scoring.
\end{itemize}

The first two are inaccessible here because no independent human annotation was
collected, so the labels produced by AlignScore \citep{zha2023alignscore} function
as the reference standard rather than as verified ground truth. Substituting an
automatic judge for human annotators is the defining move of RLAIF
\citep{bai2022constitutional, lee2023rlaif}, and it carries this limitation with
it: the resulting labels are cheap and consistent, but unaudited. The third is inaccessible because the final
study does not include a downstream policy-optimization stage
(Section~\ref{sec:methodology-evolution}). The fourth is addressed empirically and
answered in the negative: Chapter~\ref{ch:results} reports configurations in which
GCA does not lead. Claims made in later chapters are confined to learnability in
the sense defined here.

\section{Reward-Model Evaluation Protocol}
\label{sec:methodology-rm-eval}

Reward-model accuracy is evaluated by 5-fold cross-validation
\citep{kohavi1995crossvalidation} rather than a single fixed held-out split. At
the dataset sizes used here, a fixed split would leave either too little data
for training or too little for a stable accuracy estimate, whereas
cross-validation uses the full dataset for both and yields five accuracy
readings per run.

Cross-validation alone does not account for variance introduced by the reward
model's initialization and by which pairs fall into which fold. This became
apparent during the study: a first run of the final configuration produced a GCA
advantage of $+6.0$ percentage points, while a rerun of the identical
configuration \emph{at the same seed value} produced $+1.3$ percentage points
(Chapter~\ref{ch:results}, Section~\ref{sec:results-mode-sweep}). That two
same-seed runs diverge follows from an implementation detail documented in
Chapter~\ref{ch:implementation}, Section~\ref{sec:impl-rm-training}: the
\texttt{--seed} argument controls the cross-validation fold assignment, while
PyTorch's generator is left unseeded, so weight initialization, batch ordering and
dropout all vary between executions. A seed value therefore does not pin down a
run completely, and repeated runs at one seed value are not redundant.

The $n=1{,}000$ result was accordingly first pooled over \emph{six training runs}
spanning \emph{five distinct seed values} (42, 42, 7, 100, 314, 2026; seed 42
appears twice, as the original run and its rerun), giving 30 fold-level accuracy
readings per condition.

Two properties of this pooling constrain the statistical claims that can be made
from it, and are revisited in Chapter~\ref{ch:discussion},
Section~\ref{sec:disc-stats-caveat}. Six runs across five seed values is a modest
basis for a variance estimate. More consequentially, the 30 fold-level differences
are not 30 independent observations: the five folds within a run partition a
single dataset, so their training sets overlap heavily and their errors are
correlated. Procedures assuming independent observations therefore overstate the
effective sample size. Fold-level uncertainty estimates are reported in
Chapter~\ref{ch:results} as exploratory and anti-conservative, not as
confirmatory tests.

Both constraints motivated the design actually used for the reported headline
result. The \emph{run} is treated as the unit of independent replication
throughout, and the campaign was extended to \emph{twenty independently seeded
runs} (seeds 1--20, fixed before any was executed) under a corrected training
procedure in which PyTorch and CUDA are seeded per fold, so that a run is fully
determined by its seed (Chapter~\ref{ch:implementation},
Section~\ref{sec:impl-rm-training}). All confirmatory statistics reported in
this thesis --- the bootstrap interval, the Wilcoxon test, the effect sizes, and
the equivalence tests of Section~\ref{sec:results-effect-sizes} --- are computed
over these run-level differences, for which the independence assumption does
hold. Fold-level figures are retained only as descriptive summaries of spread.
Twenty runs is also the smallest campaign size that permits a two-sided Wilcoxon
result below the conventional threshold with a realistic margin: with $n$ paired
same-signed observations the smallest attainable two-sided $p$-value is
$2/2^{n}$, which at $n=6$ is $0.031$ and leaves no room for a single
discordant run, while at $n=20$ it is $1.9\times10^{-6}$
(Chapter~\ref{ch:results}, Section~\ref{sec:results-extended-campaign}).

The $n=5{,}000$ and $n=10{,}000$ configurations were initially each executed
once rather than repeated, for reasons of computational budget, and at that
stage carried correspondingly less weight than the pooled $n=1{,}000$ result.
Both were subsequently extended to twenty-run campaigns under the same
corrected, deterministic training procedure used for the $n=1{,}000$ extension,
matching its resolution exactly (Chapter~\ref{ch:results},
Section~\ref{sec:results-scale}), so all three dataset sizes are now supported
by equally powered, repeated-run evidence rather than only $n=1{,}000$.

\section{Exploratory and Confirmatory Experiments}
\label{sec:methodology-exploratory}

The experiments in Chapter~\ref{ch:results} do not all have the same evidential
status, and the distinction is material to how the results should be read.

The aggregation-formula comparison (Section~\ref{sec:results-formula-opt}), the
evaluator-mode sweep (Section~\ref{sec:results-mode-sweep}), and the
hyperparameter search (Section~\ref{sec:results-alpha0-validation}) were all
carried out on the $n=1{,}000$ setting, and the final configuration was selected
after observing their outcomes. These experiments are exploratory. The
configuration they produced was not specified in advance, and reporting it as
though it had been pre-registered would overstate the evidence: with enough
configurations examined, some separation between conditions is expected by chance
alone.

The repeated runs at $n=1{,}000$ (Section~\ref{sec:results-final-campaign}) test
whether the selected configuration produces a stable effect under repetition, but
they reuse the dataset on which that configuration was chosen and therefore cannot
establish that it generalizes beyond that setting.

The $n=5{,}000$ and $n=10{,}000$ experiments provide a weaker check than would be
ideal, for a reason that must be stated explicitly. All three subsets are drawn
from the same CNN/DailyMail test split by the same deterministic procedure at the
same seed, differing only in the number of samples requested, and they are
consequently nested rather than disjoint
(Section~\ref{sec:methodology-nesting}). The larger-scale runs therefore
re-use almost all of the data on which the configuration was selected, and are
best understood as measurements of how the effect behaves as the dataset grows,
not as out-of-sample replication on independent data. No experiment in this thesis
evaluates the selected configuration on data disjoint from the set used to choose
it.

\section{Relationship Between the Dataset Sizes}
\label{sec:methodology-nesting}

The three dataset sizes used in the final experiments are not independent samples,
and this constrains how the larger-scale runs can be interpreted.

All three subsets are produced by the same deterministic procedure
(\url{src/data/subset.py}, \texttt{select\_subset}) from the same CNN/DailyMail
test split, under the same length filters and the same subset seed (200). Only the
number of requested samples differs. Because the selection draws from one eligible
pool of 10{,}896 articles with a fixed random seed, the resulting subsets overlap
heavily. Reproducing the selection with the recorded parameters gives the
intersections in Table~\ref{tab:subset-overlap}.\footnote{Computed by re-running
\texttt{select\_subset} with the parameters recorded in
\url{configs/subset_1000.yaml} and \url{slurm/generate_candidates_scale.sh}
($\texttt{seed}=200$, test split, article limit 8{,}000 characters, summary limit
2{,}000 characters) against the local copy of the test split in
\url{data/cnn_dailymail_test_full}, and comparing the resulting SHA-256-derived
sample IDs. The candidate and preference files for these runs are not retained in
the repository, so the comparison reproduces the selection rather than reading the
stored artifacts; the selection is deterministic given these parameters.}

\begin{table}[htbp]
\centering
\caption[Overlap between the three dataset subsets]{Overlap between the subsets used in the final experiments. All three use
subset seed 200 and differ only in sample count.}
\label{tab:subset-overlap}
\begin{tabular}{lrr}
\toprule
Pair & Shared sample IDs & Share of the smaller set \\
\midrule
$n=1{,}000$ and $n=5{,}000$   & 980   & 98.0\% \\
$n=1{,}000$ and $n=10{,}000$  & 999   & 99.9\% \\
$n=5{,}000$ and $n=10{,}000$  & 5{,}000 & 100\% \\
\bottomrule
\end{tabular}
\end{table}

\begin{figure}[htbp]
\centering
\begin{tikzpicture}[font=\footnotesize]
  % eligible pool
  \draw[gray!55, fill=gray!8, rounded corners=2pt] (0,0) rectangle (11.2,4.3);
  \node[anchor=north west, gray!55!black, font=\scriptsize]
        at (0.2,4.15) {eligible pool: 10{,}896 articles};

  % n = 10000
  \draw[draw=orange!65!black, fill=orange!10, rounded corners=2pt]
        (0.45,0.35) rectangle (10.75,3.55);
  \node[anchor=north west, orange!55!black, font=\scriptsize]
        at (0.62,3.42) {$n=10{,}000$};

  % n = 5000
  \draw[draw=blue!60!black, fill=blue!10, rounded corners=2pt]
        (0.85,0.7) rectangle (5.95,2.95);
  \node[anchor=north west, blue!50!black, font=\scriptsize]
        at (1.02,2.82) {$n=5{,}000$};

  % n = 1000
  \draw[draw=black!70, fill=white, rounded corners=2pt]
        (1.25,1.05) rectangle (2.85,2.25);
  \node[anchor=north west, font=\scriptsize] at (1.4,2.12) {$n=1{,}000$};

  % annotations with leader lines, placed in free space
  \node[anchor=west, font=\scriptsize, black!75] (a1) at (6.35,2.55)
       {980 of 1{,}000 also in $n{=}5{,}000$};
  \draw[black!45, -{Stealth[length=1.5mm]}] (a1.west) -- (2.95,2.0);

  \node[anchor=west, font=\scriptsize, black!75] (a2) at (6.35,1.85)
       {999 of 1{,}000 also in $n{=}10{,}000$};
  \draw[black!45, -{Stealth[length=1.5mm]}] (a2.west) -- (2.95,1.55);

  \node[anchor=west, font=\scriptsize, blue!50!black] (a3) at (6.35,1.15)
       {all 5{,}000 also in $n{=}10{,}000$};
  \draw[blue!40, -{Stealth[length=1.5mm]}] (a3.west) -- (6.0,1.0);
\end{tikzpicture}
\caption[Nesting of the three dataset subsets]{The three subsets are nested rather than independent. Box sizes are
schematic; the counts are the measured intersections reported in
Table~\ref{tab:subset-overlap}.}
\label{fig:subset-nesting}
\end{figure}

The $n=5{,}000$ subset is contained entirely within the $n=10{,}000$ subset, and
the $n=1{,}000$ subset is almost entirely contained within both. The
$n=10{,}000$ set additionally covers 91.8\% of the whole eligible pool, so at that
size little of the test split remains unused in any case.

Two consequences follow, and both are carried through the rest of the thesis. The
larger-scale experiments do not constitute out-of-sample replication: they re-use
almost all of the data on which the configuration was selected, and therefore
measure how the effect behaves as further data is added to the same pool rather
than whether it transfers to new data. And because the samples are nested, the
three results are not statistically independent of one another, so the differences
between them should not be read as three separate estimates of the same quantity.
Drawing the larger subsets with different seeds would have made them independent,
at the cost of the nesting property the generation script was written to provide;
this is recorded as a limitation in Chapter~\ref{ch:conclusion}.

\section{Independent Ground-Truth Construction}
\label{sec:methodology-ground-truth}

The evaluation described above has two properties that limit what it can
establish, both noted in Section~\ref{sec:methodology-exploratory} and
Section~\ref{sec:concl-limitations}: the labels a trained reward model is
scored against are supplied by the same evaluator that constructed its
training preferences, and the $n=5{,}000$/$n=10{,}000$ checks reuse almost
all of the data the configuration was selected on rather than evaluating it
on disjoint data. RQ4 was formulated to address both at once, by scoring the
already-trained reward models against a held-out set built specifically to
be disjoint from every training pool used in this thesis, under an
evaluation criterion that does not depend on how any individual training
pair happened to be labelled.

\paragraph{Held-out article selection.} 500 articles were drawn from the same
CNN/DailyMail test split, using \texttt{select\_disjoint\_subset}
(\url{src/data/subset.py}): the function reconstructs the exact
\texttt{rng.sample()} call that produced the $n=10{,}000$, seed-200 draw used
throughout this thesis, removes those indices from the eligible pool, and
draws 500 samples from what remains under a new seed (999). Because the
$n=5{,}000$ subset is fully contained in the $n=10{,}000$ draw and the
$n=1{,}000$ subset is 99.9\% contained in it
(Table~\ref{tab:subset-overlap}), excluding the $n=10{,}000$ draw is
sufficient to make the held-out set disjoint from all three training pools
used in the main campaign; two further training pools introduced later in
this work, at $n=2{,}000$ and $n=3{,}000$ (Section~\ref{sec:results-scale-curve}),
were verified by direct index comparison to be exact subsets of the
$n=10{,}000$ draw and are therefore disjoint from the held-out set for the
same reason.

\paragraph{Candidate generation and scoring.} Two candidate summaries were
generated per held-out article with the same Mistral-7B-Instruct-v0.3 pipeline
and the same two decoding temperatures used throughout this thesis
(Section~\ref{sec:methodology-design}), giving 1{,}000 candidate summaries in
total. Every candidate was scored with AlignScore under the locked
configuration used for the main campaign: sentence-level scoring, GCA
aggregation with $\alpha=0.0$ (Section~\ref{sec:results-formula-opt}). This
GCA-aggregated score is used only to construct the evaluation criterion
below; it is never used as a training signal for the held-out set, and no
reward model is trained on these 500 articles.

\paragraph{Evaluation criterion.} Two ground-truth constructions were used.
The first pairs each article's own two candidates directly, mirroring the
structure of the training preferences; the two candidates are generated by
the same model at similar temperatures, so their quality gap is typically
small (mean GCA score gap approximately 0.04). The second pools all
1{,}000 scored candidates across all 500 articles, sorts them by GCA score,
and compares the top $X\%$ against the bottom $X\%$ for $X\in\{25,10,5\}$,
which produces progressively larger, more unambiguous quality gaps between
the two sides at the cost of fewer summaries per comparison
(Section~\ref{sec:results-ground-truth}). At $X=5$, every high-scoring
summary is compared against every low-scoring summary (2{,}500 pairs from 50
summaries per side) rather than one random pairing, to remove pairing-order
noise from the smallest, highest-signal condition. A trained reward model's
\emph{success rate} on a given construction is the fraction of comparisons in
which it scores the higher-GCA-score summary above the lower-GCA-score one;
this is the dependent variable throughout Section~\ref{sec:results-ground-truth},
and it is defined identically for the holistic and GCA reward models, so
neither construction favours either condition by design.

\paragraph{Checkpoint retention.} The training procedure used for RQ1--RQ3
(Section~\ref{sec:methodology-rm-eval}) discards each fold's trained weights
after recording its accuracy, since only the accuracy is needed there. RQ4
requires scoring the same trained model against a second, independent
criterion, so the training procedure was changed for this evaluation to save
the full reward-model checkpoint (encoder weights and reward head) for every
run, rather than only its cross-validation accuracy.

\section{Evolution of the Study Design}
\label{sec:methodology-evolution}

The study reported here is narrower in scope than the one outlined in the original
research plan, and the reasons are methodological.

The initial design used the generated preferences directly for DPO-based policy
optimization, evaluating factual consistency on the fine-tuned model's own
generations. Early downstream experiments made it difficult to attribute observed
differences to any single stage: a change in generated-summary factuality could
originate in the preference-construction procedure, in the quality of the reward
signal, or in the policy-optimization step, and the design provided no way to
separate these. Increasing the number of stages also increased the number of
variance sources standing between the manipulated variable and the measured
outcome.

The final study therefore isolates an earlier and more controlled diagnostic
question: whether the preference data produced by each construction procedure is
recoverable by a fixed reward-model architecture. Removing the policy-optimization
stage removes both reinforcement-learning dynamics and decoding variance from the
measurement path, leaving a comparison in which granularity is the principal
difference between conditions.

The claims of this thesis are correspondingly narrower. The work evaluates the
learnability of the preference signal rather than demonstrating improvements in
generated-summary factuality (Section~\ref{sec:methodology-scope}). Two further
consequences follow from the same change. Margin-based filtering of near-tie pairs
was removed, since AlignScore produces continuous scores for which exact ties are
effectively impossible, and the threshold discarded roughly one fifth of the data
without a demonstrated benefit. The reliability protocol designed for a generative
LLM judge, comprising confidence gating, order randomization, and rationale
requirements, no longer applies once the evaluator is a deterministic model that
is never prompted and returns no rationale; a planned human audit of judge
rationales was not carried out and is recorded as a limitation in
Chapter~\ref{ch:conclusion}.

\section{Validity Considerations}
\label{sec:methodology-ethics}

\paragraph{Ethical and practical considerations.} All experiments use
CNN/DailyMail \citep{hermann2015teaching, nallapati2016abstractive}, a publicly
available corpus of news articles and human-written summaries; no personally
sensitive data is collected. Because both the evaluator's
judgments and the generated summaries can be incorrect, results are reported as
measurements of preference-data learnability rather than as statements about
factual correctness, and model outputs are not presented as verified fact. Compute
use was bounded by working with fixed-size subsets.

\paragraph{Candidate generation and decoding temperature.} The two candidates in
each pair are produced by the same model at different decoding temperatures
($T=0.7$ and $T=1.0$). This introduces a potential confound: preference structure
may correlate with systematic stylistic differences induced by temperature rather
than with factuality alone. The preference data stored in the repository allow the
asymmetry to be quantified directly, and Table~\ref{tab:temperature-asymmetry}
reports it.\footnote{Computed from
\url{data/preferences/holistic_reward_preferences_200.jsonl} and
\url{data/preferences/gca_reward_preferences_200.jsonl}. The margin-0 rows
recompute the decision from the stored scores at the threshold used in the final
configuration; the margin-0.05 rows reproduce the stored decisions, whose
aggregate counts also appear in
\url{reports/reward_model_judging_summary.csv}. Preference files for the final
$n=1{,}000$ configuration are not retained in the repository, so these proportions
could not be recomputed under the final evaluator mode.}

\begin{table}[htbp]
\centering
\caption[Decoding-temperature asymmetry in the preference labels]{Proportion of pairs in which the lower-temperature candidate ($T=0.7$)
is preferred, on the $n=200$ preference set (evaluator mode \texttt{nli\_sp},
$\alpha=0.5$).}
\label{tab:temperature-asymmetry}
\begin{tabular}{llrrr}
\toprule
Condition & Threshold & $T{=}0.7$ preferred & $T{=}1.0$ preferred & Share \\
\midrule
Holistic & margin $=0$    & 137 & 63 & 68.5\% \\
GCA      & margin $=0$    & 126 & 74 & 63.0\% \\
\addlinespace
Holistic & margin $=0.05$ & 109 & 45 & 70.8\% \\
GCA      & margin $=0.05$ &  97 & 60 & 61.8\% \\
\bottomrule
\end{tabular}
\end{table}

The asymmetry is substantial under both conditions and consistently larger under
holistic scoring, and it is not an artifact of the margin threshold. Two
implications follow. Features correlated with decoding temperature, such as
summary length or lexical conservativeness, may carry a non-trivial share of the
label information, so a reward model could recover part of the preference relation
without representing factuality as such. And because the asymmetry differs between
conditions by roughly five percentage points, the two preference sets are not
equivalent in this respect either, which means the comparison between them is not
perfectly clean with respect to temperature.

These figures come from the $n=200$ set scored under the earlier evaluator
configuration and need not hold exactly at the final one. They are reported as
evidence that the confound is real and measurable, not as a quantification of its
effect on the reported results. Controlling candidate generation more explicitly,
by sampling both candidates at one temperature, would remove the confound at
source and is left to future work (Chapter~\ref{ch:conclusion}).

\paragraph{How much label information do surface features actually carry?} The
concern above is only as serious as the amount of label information that
temperature-correlated surface features actually carry, and that is measurable
directly on the same stored data. Table~\ref{tab:surface-baselines} reports the
accuracy of the cheapest possible classifiers --- ones that ignore the article
entirely and decide from a single surface property of the two candidate
summaries --- against the same decisions the evaluator
made.\footnote{Computed by \url{analysis/surface_feature_baselines.py} from the
same two preference files as Table~\ref{tab:temperature-asymmetry}, restricted
to the pairs carrying an A/B decision (154 holistic, 157 GCA).}

\begin{table}[htbp]
\centering
\caption[Surface-feature baselines against the judge's decisions]{Accuracy of surface-feature classifiers against the evaluator's own
decisions, $n=200$ preference set (mode \texttt{nli\_sp}, $\alpha=0.5$). Each
row states the accuracy of the better of the two opposing rules (prefer
longer/shorter, prefer more/fewer sentences). Mean lengths are in characters.}
\label{tab:surface-baselines}
\begin{tabular}{lrr}
\toprule
Predictor & Holistic & GCA \\
\midrule
Prefer the longer summary        & 0.494 & 0.465 \\
Prefer the shorter summary       & 0.506 & 0.535 \\
Prefer more sentences            & 0.519 & 0.475 \\
Prefer fewer sentences           & 0.481 & 0.525 \\
\addlinespace
Mean length, chosen              & 865.6 & 840.5 \\
Mean length, rejected            & 863.5 & 846.2 \\
\bottomrule
\end{tabular}
\end{table}

None of these predictors performs meaningfully above chance. The mean length of
the chosen summary differs from that of the rejected summary by 2 characters
under holistic scoring and by $-6$ characters under GCA, on summaries averaging
roughly 850 characters, and sentence counts differ by 0.01 and $-0.28$
respectively. Length and sentence count, the two most obvious surface proxies
for decoding temperature, therefore do not carry the preference signal, and a
reward model cannot be recovering its accuracy from them.

This narrows the confound without eliminating it, and the distinction matters.
The 68.5\% and 63.0\% asymmetries in
Table~\ref{tab:temperature-asymmetry} are not themselves a shortcut available to
the reward model: the model is trained on (article, chosen) and (article,
rejected) text pairs and never observes which candidate came from which
temperature, or which position it occupied
(Chapter~\ref{ch:implementation}, Section~\ref{sec:impl-rm-training}). For the
asymmetry to inflate reward-model accuracy, some \emph{textual} property
correlated with temperature would have to be learnable from the summary itself.
Table~\ref{tab:surface-baselines} rules out the two most likely candidates for
such a property. Subtler lexical or stylistic correlates are not ruled out, and
testing them would require the feature-level analysis recorded as future work,
but the simple version of the objection does not hold.

\paragraph{Threats to validity.} Table~\ref{tab:threats-to-validity} summarizes
the threats considered and how each is handled.

\begin{table}[htbp]
\centering
\caption[Threats to validity]{Threats to validity and how each is addressed.}
\label{tab:threats-to-validity}
\begin{tabularx}{\textwidth}{lX}
\toprule
Threat & Status \\
\midrule
Single evaluator (AlignScore) & Not mitigated; all results are conditional on this
  evaluator's judgments, which are not independent ground truth. \\
Evaluator and aggregation sensitivity & Measured directly via the formula
  comparison and evaluator-mode sweep (Chapter~\ref{ch:results}). \\
Fold- and run-level noise & Addressed by twenty independently seeded runs at each
  of the three dataset sizes, under fully deterministic per-fold seeding
  (Section~\ref{sec:methodology-rm-eval}). \\
Non-independence of pooled folds & Mitigated by testing at the run level, where
  the independence assumption holds; fold-level figures are reported only as
  descriptive summaries of spread (Section~\ref{sec:methodology-rm-eval}). \\
Configuration selected on observed results & Acknowledged, not remedied; the
  larger-scale subsets are nested supersets of the selection data
  (Section~\ref{sec:methodology-nesting}), so no run evaluates the configuration
  on independent data. \\
Scale sensitivity & Examined explicitly rather than assumed away
  (Section~\ref{sec:results-scale}). \\
Decoding-temperature confound & Quantified where data permit; not controlled for
  in the design. Length and sentence count are excluded as carriers of the
  signal (Table~\ref{tab:surface-baselines}); subtler lexical correlates remain
  untested (see above). \\
Source-document truncation & Documented in Chapter~\ref{ch:implementation},
  Section~\ref{sec:impl-rm-training}; discussed as a limitation. \\
Weak absolute accuracy & Both conditions remain near chance
  (Chapter~\ref{ch:discussion}, Section~\ref{sec:disc-weak-signal}). \\
\bottomrule
\end{tabularx}
\end{table}
